\documentclass[a4paper]{article}

\usepackage{graphicx}
\usepackage{amsmath,amssymb}
\usepackage{minted}
\usepackage{multirow}
\usepackage{float}
\usepackage{todonotes}
\usepackage{forest}
\usepackage{xstring}
\usepackage{xcolor}
\usepackage[most]{tcolorbox}
\definecolor{TranslationGray}{RGB}{120,120,120}
\newcommand{\EN}[1]{\par{\small\textcolor{TranslationGray}{#1}}\vspace{0.25em}}
\usepackage{xparse}
\usepackage{natbib}

\usetikzlibrary{graphs,graphdrawing}
\usegdlibrary{layered}

% for external graphviz
\input{/home/donkarlo/Dropbox/repo/nd_ai_project/src/nd_ai/knoweldge_representation/language/formal/computer/programming/tex/latex/code_clips/external_graphviz.tex}

% for tags
\input{/home/donkarlo/Dropbox/repo/nd_ai_project/src/nd_ai/knoweldge_representation/language/formal/computer/programming/tex/latex/parts/control_sequence/kind/semtag/semtag.tex}

% for links
\input{/home/donkarlo/Dropbox/repo/nd_ai_project/src/nd_ai/knoweldge_representation/language/formal/computer/programming/tex/latex/code_clips/seehere.tex}

\title{Kausale Mittel im Deutschen}
\date{}

\begin{document}
    \maketitle
    
    
    \section{Einleitung}
        Kausale Mittel drücken im Deutschen einen \textbf{Grund}, eine \textbf{Ursache}, einen \textbf{Anlass} oder eine \textbf{Folge} aus. Sie helfen dabei, Zusammenhänge zwischen Ereignissen, Handlungen und Situationen sprachlich klar zu machen. Sehr oft antworten sie auf Fragen wie \textbf{Warum?}, \textbf{Weshalb?}, \textbf{Wieso?} oder \textbf{Aus welchem Grund?}
        
        Man kann kausale Mittel in verschiedene Gruppen einteilen. Einige nennen direkt die Ursache, andere drücken die Folge aus. Manche erscheinen als Präpositionen, andere als Konjunktionen, Adverbien, Verben oder Nomen. Für das Verständnis ist besonders wichtig, ob ein Ausdruck eher den \textbf{Grund} oder eher die \textbf{Folge} bezeichnet.
    
    
    \section{Grundbegriffe}
        
        \subsection{Grund, Ursache, Anlass, Folge}
            \textbf{der Grund} bezeichnet allgemein den Anlass oder den sachlichen Hintergrund für etwas.\\
            \EN{the reason}
            
            \textbf{die Ursache} ist oft etwas sachlicher und bezeichnet den Auslöser eines Ereignisses.\\
            \EN{the cause}
            
            \textbf{der Anlass} bezeichnet häufig einen konkreten Auslöser oder einen äußeren Grund für eine Handlung.\\
            \EN{the occasion; the reason}
            
            \textbf{die Folge} ist das Resultat oder die Konsequenz eines Grundes oder einer Ursache.\\
            \EN{the consequence; the result}
        
        \subsection{Kausaler Zusammenhang}
            Ein \textbf{kausaler Zusammenhang} liegt vor, wenn zwischen zwei Aussagen eine Beziehung von Ursache und Folge besteht. Eine Sache geschieht, \textbf{weil} etwas anderes der Grund dafür ist, oder etwas passiert und \textbf{deshalb} ergibt sich daraus eine Konsequenz.
            
            \paragraph{Beispiele}
                Es regnet, deshalb bleiben wir zu Hause.\\
                \EN{It is raining, therefore we are staying at home.}
                
                Wir bleiben zu Hause, weil es regnet.\\
                \EN{We are staying at home because it is raining.}
    
    
    \section{Kausale Präpositionen}
    Kausale Präpositionen verbinden meist ein Nomen oder eine Nominalgruppe mit dem Rest des Satzes und geben einen Grund oder eine Ursache an.
    
    \subsection{wegen + Genitiv}
        \textbf{wegen} nennt einen Grund oder eine Ursache. Im heutigen Sprachgebrauch kommt in der Umgangssprache manchmal auch der Dativ vor, aber standardsprachlich gilt vor allem der \textbf{Genitiv}.
        
        \paragraph{Beispiele}
            Wegen des Regens bleiben wir zu Hause.\\
            \EN{Because of the rain, we are staying at home.}
            
            Wegen meiner Krankheit konnte ich nicht kommen.\\
            \EN{Because of my illness, I could not come.}
            
            Wegen eines Missverständnisses gab es ein Problem.\\
            \EN{Because of a misunderstanding, there was a problem.}
    
    \subsection{aufgrund + Genitiv}
    \textbf{aufgrund} bedeutet fast dasselbe wie \textbf{wegen}, wirkt aber meist etwas formeller und wird besonders in schriftlichen oder offiziellen Kontexten verwendet.
    
    \paragraph{Beispiele}
        Aufgrund eines Fehlers wurde der Termin verschoben.\\
        \EN{Because of an error, the appointment was postponed.}
        
        Aufgrund der Wetterlage blieb die Schule geschlossen.\\
        \EN{Because of the weather situation, the school remained closed.}
        
        Aufgrund neuer Informationen wurde die Entscheidung geändert.\\
        \EN{Because of new information, the decision was changed.}
    
    \subsection{infolge + Genitiv}
    \textbf{infolge} ist sehr formell und bezeichnet eine Folge, deren Ursache im Genitiv genannt wird. Semantisch liegt der Fokus oft auf dem eingetretenen Resultat.
    
    \paragraph{Beispiele}
        Infolge des Unfalls gab es eine Straßensperre.\\
        \EN{As a result of the accident, there was a road closure.}
        
        Infolge der Krise verloren viele Menschen ihre Arbeit.\\
        \EN{As a result of the crisis, many people lost their jobs.}
        
        Infolge des starken Schneefalls kam es zu Verspätungen.\\
        \EN{As a result of heavy snowfall, there were delays.}
    
    \subsection{aus + Dativ}
    \textbf{aus} drückt oft einen inneren Beweggrund, ein Motiv oder eine Haltung aus. Es geht häufig nicht um eine äußere sachliche Ursache, sondern um eine persönliche Motivation.
    
    \paragraph{Beispiele}
        Er half ihr aus Freundschaft.\\
        \EN{He helped her out of friendship.}
        
        Sie tat das aus Liebe.\\
        \EN{She did that out of love.}
        
        Aus Angst sagte er nichts.\\
        \EN{Out of fear, he said nothing.}
    
    \subsection{vor + Dativ}
    \textbf{vor} wird bei starken Gefühlen oder körperlichen Reaktionen gebraucht. Es bezeichnet oft einen emotionalen oder körperlichen Grund.
    
    \paragraph{Beispiele}
        Sie zitterte vor Angst.\\
        \EN{She was trembling with fear.}
        
        Er weinte vor Freude.\\
        \EN{He cried with joy.}
        
        Das Kind konnte vor Aufregung nicht schlafen.\\
        \EN{The child could not sleep because of excitement.}
    
    \subsection{durch + Akkusativ}
    \textbf{durch} kann Ursache, Mittel oder auslösenden Faktor bezeichnen. Im kausalen Gebrauch steht oft der Faktor im Vordergrund, der etwas bewirkt.
    
    \paragraph{Beispiele}
        Durch den Lärm konnte ich nicht schlafen.\\
        \EN{Because of the noise, I could not sleep.}
        
        Durch seine Hilfe war alles einfacher.\\
        \EN{Because of his help, everything was easier.}
        
        Durch einen technischen Fehler fiel das System aus.\\
        \EN{Because of a technical error, the system failed.}
    
    
    \section{Kausale Konjunktionen}
    Kausale Konjunktionen verbinden Sätze miteinander und leiten typischerweise einen Nebensatz mit Grund ein.
    
    \subsection{weil}
        \textbf{weil} ist die häufigste kausale Konjunktion. Sie leitet einen Nebensatz ein, in dem das konjugierte Verb am Ende steht. \textbf{weil} nennt direkt den Grund.
        
        \paragraph{Beispiele}
            Ich bleibe zu Hause, weil ich krank bin.\\
            \EN{I am staying at home because I am ill.}
            
            Er lernt Deutsch, weil er in Österreich arbeitet.\\
            \EN{He learns German because he works in Austria.}
            
            Wir kommen später, weil der Zug Verspätung hat.\\
            \EN{We are coming later because the train is delayed.}
    
    \subsection{da}
    \textbf{da} ist \textbf{weil} sehr ähnlich, klingt aber oft etwas formeller. Häufig wird \textbf{da} verwendet, wenn der Grund als bekannt, selbstverständlich oder bereits eingeführt erscheint.
    
    \paragraph{Beispiele}
        Da es spät ist, gehen wir jetzt.\\
        \EN{Since it is late, we are leaving now.}
        
        Da ich morgen früh aufstehen muss, gehe ich schlafen.\\
        \EN{Since I have to get up early tomorrow, I am going to sleep.}
        
        Da er keine Zeit hatte, sagte er ab.\\
        \EN{Since he had no time, he cancelled.}
    
    \subsection{zumal}
    \textbf{zumal} fügt einen besonders wichtigen oder zusätzlichen Grund hinzu. Es hebt oft einen Grund hervor, der die Aussage noch stärker stützt.
    
    \paragraph{Beispiele}
        Wir bleiben hier, zumal das Wetter schlecht ist.\\
        \EN{We are staying here, especially since the weather is bad.}
        
        Er sollte sich ausruhen, zumal er noch Fieber hat.\\
        \EN{He should rest, especially since he still has a fever.}
        
        Die Reise war anstrengend, zumal wir kaum geschlafen hatten.\\
        \EN{The trip was exhausting, especially since we had hardly slept.}
    
    
    \section{Kausale Adverbien und Adverbialausdrücke}
    Diese Ausdrücke beziehen sich meist auf einen vorher genannten Grund und formulieren daraus eine Folge oder Konsequenz.
    
    \subsection{deshalb}
        \textbf{deshalb} bedeutet \textbf{aus diesem Grund}. Es bezieht sich auf eine vorher genannte Ursache und leitet typischerweise die Folge ein.
        
        \paragraph{Beispiele}
            Ich war müde, deshalb bin ich früh ins Bett gegangen.\\
            \EN{I was tired, therefore I went to bed early.}
            
            Es regnet, deshalb nehme ich einen Schirm mit.\\
            \EN{It is raining, therefore I am taking an umbrella.}
            
            Er hat viel gearbeitet, deshalb braucht er eine Pause.\\
            \EN{He worked a lot, therefore he needs a break.}
    
    \subsection{deswegen}
    \textbf{deswegen} ist in der Bedeutung sehr nah bei \textbf{deshalb}. Es ist ein kausales Adverb und drückt eine Folge aus, die sich aus einem zuvor genannten Grund ergibt.
    
    \paragraph{Beispiele}
        Es regnet, deswegen nehme ich einen Schirm mit.\\
        \EN{It is raining, therefore I am taking an umbrella.}
        
        Ich war krank, deswegen konnte ich nicht kommen.\\
        \EN{I was ill, therefore I could not come.}
        
        Sie hatte wenig Zeit, deswegen arbeitete sie sehr schnell.\\
        \EN{She had little time, therefore she worked very quickly.}
    
    \subsection{daher}
    \textbf{daher} bedeutet ebenfalls \textbf{aus diesem Grund}. Es wirkt oft etwas schriftsprachlicher oder nüchterner.
    
    \paragraph{Beispiele}
        Er hat viel gearbeitet, daher ist er erschöpft.\\
        \EN{He worked a lot, therefore he is exhausted.}
        
        Die Daten sind unvollständig, daher ist das Ergebnis unsicher.\\
        \EN{The data are incomplete, therefore the result is uncertain.}
        
        Es bestand ein Risiko, daher wurden zusätzliche Maßnahmen getroffen.\\
        \EN{There was a risk, therefore additional measures were taken.}
    
    \subsection{darum}
    \textbf{darum} ist ebenfalls ein Ausdruck der Folge. Im Vergleich zu \textbf{daher} wirkt es oft etwas neutraler oder gesprochener.
    
    \paragraph{Beispiele}
        Ich habe keine Zeit, darum komme ich nicht mit.\\
        \EN{I do not have time, therefore I am not coming along.}
        
        Sie war müde, darum ging sie nach Hause.\\
        \EN{She was tired, therefore she went home.}
        
        Er hat lange geübt, darum spielt er jetzt besser.\\
        \EN{He practised for a long time, therefore he plays better now.}
    
    \subsection{aus diesem Grund}
    \textbf{aus diesem Grund} ist ein expliziter, klarer Ausdruck für eine Folgerung. Er ist besonders in schriftlichen Erklärungen und argumentativen Texten nützlich.
    
    \paragraph{Beispiele}
        Aus diesem Grund müssen wir vorsichtig sein.\\
        \EN{For this reason, we must be careful.}
        
        Aus diesem Grund wurde die Regel geändert.\\
        \EN{For this reason, the rule was changed.}
        
        Aus diesem Grund kann der Plan nicht sofort umgesetzt werden.\\
        \EN{For this reason, the plan cannot be implemented immediately.}
    
    \subsection{aus diesem Anlass}
    \textbf{aus diesem Anlass} bezieht sich stärker auf einen konkreten Anlass oder eine Gelegenheit als auf eine strenge Ursache im logischen Sinn.
    
    \paragraph{Beispiele}
        Aus diesem Anlass wurde eine Feier organisiert.\\
        \EN{For this occasion, a celebration was organised.}
        
        Aus diesem Anlass hielt der Direktor eine Rede.\\
        \EN{On this occasion, the director gave a speech.}
        
        Aus diesem Anlass traf sich die ganze Familie.\\
        \EN{On this occasion, the whole family met.}
    
    
    \section{Kausale Fragewörter}
    Kausale Fragewörter dienen dazu, nach dem Grund, der Ursache oder dem Anlass zu fragen.
    
    \subsection{warum}
        \textbf{warum} ist das häufigste Fragewort für Gründe.
        
        \paragraph{Beispiele}
            Warum lernst du Deutsch?\\
            \EN{Why are you learning German?}
            
            Warum bist du zu spät gekommen?\\
            \EN{Why did you come late?}
            
            Warum hat er so reagiert?\\
            \EN{Why did he react like that?}
    
    \subsection{weshalb}
    \textbf{weshalb} ist \textbf{warum} sehr ähnlich, wirkt aber manchmal etwas formeller.
    
    \paragraph{Beispiele}
        Weshalb bist du zu spät gekommen?\\
        \EN{Why did you come late?}
        
        Weshalb wurde das Treffen verschoben?\\
        \EN{Why was the meeting postponed?}
        
        Weshalb möchte sie nicht mitkommen?\\
        \EN{Why does she not want to come along?}
    
    \subsection{wieso}
    \textbf{wieso} ist im gesprochenen Deutsch sehr häufig und wirkt oft etwas alltagssprachlicher.
    
    \paragraph{Beispiele}
        Wieso lachst du?\\
        \EN{Why are you laughing?}
        
        Wieso bist du so still?\\
        \EN{Why are you so quiet?}
        
        Wieso funktioniert das nicht?\\
        \EN{Why does that not work?}
    
    \subsection{aus welchem Grund}
    \textbf{aus welchem Grund} ist eine formellere und explizitere Art, nach der Ursache zu fragen.
    
    \paragraph{Beispiele}
        Aus welchem Grund wurde der Kurs abgesagt?\\
        \EN{For what reason was the course cancelled?}
        
        Aus welchem Grund lehnen Sie den Vorschlag ab?\\
        \EN{For what reason do you reject the proposal?}
        
        Aus welchem Grund hat sich die Entscheidung geändert?\\
        \EN{For what reason has the decision changed?}
    
    
    \section{Verben mit kausaler Bedeutung}
    Auch manche Verben tragen selbst schon eine kausale Bedeutung in sich. Sie bezeichnen, dass etwas etwas anderes auslöst, verursacht oder zur Folge hat.
    
    \subsection{verursachen}
        \textbf{verursachen} bedeutet, die Ursache für etwas zu sein.
        
        \paragraph{Beispiele}
            Der Sturm verursachte große Schäden.\\
            \EN{The storm caused major damage.}
            
            Das Missverständnis verursachte einen Konflikt.\\
            \EN{The misunderstanding caused a conflict.}
            
            Zu wenig Schlaf kann Konzentrationsprobleme verursachen.\\
            \EN{Too little sleep can cause concentration problems.}
    
    \subsection{bewirken}
    \textbf{bewirken} bedeutet, etwas hervorzurufen oder eine Wirkung zu erzeugen.
    
    \paragraph{Beispiele}
        Die Übung bewirkt eine Verbesserung.\\
        \EN{The exercise causes an improvement.}
        
        Die Reform bewirkte viele Veränderungen.\\
        \EN{The reform brought about many changes.}
        
        Dieses Medikament bewirkt oft schnelle Linderung.\\
        \EN{This medicine often brings quick relief.}
    
    \subsection{führen zu + Dativ}
    \textbf{führen zu} bezeichnet, dass eine Sache eine bestimmte Folge nach sich zieht.
    
    \paragraph{Beispiele}
        Der Fehler führte zu einem Problem.\\
        \EN{The error led to a problem.}
        
        Zu viel Stress führt zu Erschöpfung.\\
        \EN{Too much stress leads to exhaustion.}
        
        Die Diskussion führte zu keiner Lösung.\\
        \EN{The discussion led to no solution.}
    
    \subsection{auslösen}
    \textbf{auslösen} bedeutet, etwas in Gang zu setzen oder hervorzurufen.
    
    \paragraph{Beispiele}
        Die Nachricht löste Panik aus.\\
        \EN{The news triggered panic.}
        
        Das Foto löste bei ihm Erinnerungen aus.\\
        \EN{The photo triggered memories in him.}
        
        Die Kritik löste eine Debatte aus.\\
        \EN{The criticism triggered a debate.}
    
    \subsection{hervorrufen}
    \textbf{hervorrufen} ist ähnlich wie \textbf{auslösen}, wirkt aber oft etwas schriftsprachlicher.
    
    \paragraph{Beispiele}
        Das Medikament kann Müdigkeit hervorrufen.\\
        \EN{The medication can cause tiredness.}
        
        Sein Verhalten rief Misstrauen hervor.\\
        \EN{His behaviour caused mistrust.}
        
        Das Ereignis rief starke Reaktionen hervor.\\
        \EN{The event caused strong reactions.}
    
    
    \section{Nominalstil mit kausaler Bedeutung}
    Im Deutschen kann man kausale Zusammenhänge nicht nur mit Präpositionen, Konjunktionen oder Adverbien ausdrücken, sondern auch im \textbf{Nominalstil}. Das ist besonders in schriftlichen, wissenschaftlichen oder formellen Texten häufig.
    
    \subsection{der Grund}
        
        \paragraph{Beispiele}
            Der Grund für seine Abwesenheit ist unklar.\\
            \EN{The reason for his absence is unclear.}
            
            Es gibt einen guten Grund für diese Entscheidung.\\
            \EN{There is a good reason for this decision.}
            
            Den Grund kenne ich nicht.\\
            \EN{I do not know the reason.}
    
    \subsection{die Ursache}
    
    \paragraph{Beispiele}
        Die Ursache des Problems wurde gefunden.\\
        \EN{The cause of the problem was found.}
        
        Die genaue Ursache ist noch unbekannt.\\
        \EN{The exact cause is still unknown.}
        
        Man sucht nach der Ursache des Fehlers.\\
        \EN{They are looking for the cause of the error.}
    
    \subsection{der Anlass}
    
    \paragraph{Beispiele}
        Der Anlass für das Treffen war ein Geburtstag.\\
        \EN{The occasion for the meeting was a birthday.}
        
        Es gab keinen Anlass zur Sorge.\\
        \EN{There was no reason for concern.}
        
        Was war der Anlass für seine Reise?\\
        \EN{What was the reason for his trip?}
    
    \subsection{der Auslöser}
    
    \paragraph{Beispiele}
        Der Auslöser des Streits war eine Kleinigkeit.\\
        \EN{The trigger of the argument was a small thing.}
        
        Ein Kommentar war der Auslöser der Debatte.\\
        \EN{One comment was the trigger of the debate.}
        
        Man kennt den Auslöser noch nicht genau.\\
        \EN{The trigger is not yet known exactly.}
    
    
    \section{Nützliche Unterscheidungen}
    
    \subsection{wegen und aufgrund}
        \textbf{wegen} und \textbf{aufgrund} sind in der Bedeutung oft ähnlich. \textbf{wegen} ist neutral und sehr häufig. \textbf{aufgrund} wirkt meist formeller und schriftsprachlicher.
        
        \paragraph{Beispiele}
            Wegen des Wetters bleiben wir zu Hause.\\
            \EN{Because of the weather, we are staying at home.}
            
            Aufgrund des Wetters bleibt die Veranstaltung geschlossen.\\
            \EN{Because of the weather, the event remains closed.}
    
    \subsection{aus und wegen}
    \textbf{aus} drückt häufig ein inneres Motiv aus, \textbf{wegen} eher einen allgemeinen äußeren oder neutralen Grund.
    
    \paragraph{Beispiele}
        Aus Liebe blieb sie bei ihm.\\
        \EN{She stayed with him out of love.}
        
        Wegen des Wetters blieb sie zu Hause.\\
        \EN{Because of the weather, she stayed at home.}
    
    \subsection{vor und aus}
    \textbf{vor} wird besonders bei starken emotionalen oder körperlichen Reaktionen verwendet, \textbf{aus} eher bei bewussten inneren Motiven.
    
    \paragraph{Beispiele}
        Er zitterte vor Angst.\\
        \EN{He trembled with fear.}
        
        Er schwieg aus Angst.\\
        \EN{He kept silent out of fear.}
    
    \subsection{weil und deshalb}
    \textbf{weil} nennt den Grund direkt im Nebensatz. \textbf{deshalb} nimmt einen vorher genannten Grund auf und formuliert daraus die Folge.
    
    \paragraph{Beispiele}
        Ich gehe nach Hause, weil ich müde bin.\\
        \EN{I am going home because I am tired.}
        
        Ich bin müde, deshalb gehe ich nach Hause.\\
        \EN{I am tired, therefore I am going home.}
    
    \subsection{wegen und deswegen}
    Diese beiden Ausdrücke werden oft zusammen gelernt, gehören aber semantisch nicht genau zur gleichen Untergruppe.
    
    \textbf{wegen} bezeichnet den \textbf{Grund} oder die \textbf{Ursache}.\\
    \textbf{deswegen} bezeichnet die \textbf{Folge} oder \textbf{Konsequenz}.
    
    \paragraph{Beispiele}
        Wegen des Regens bleiben wir zu Hause.\\
        \EN{Because of the rain, we are staying at home.}
        
        Es regnet, deswegen bleiben wir zu Hause.\\
        \EN{It is raining, therefore we are staying at home.}
    
    
    \section{Semantische Einordnung}
    Wenn man die Ausdrücke \textbf{wegen}, \textbf{deswegen}, \textbf{weil}, \textbf{darum}, \textbf{daher} und ähnliche Wörter nicht nach ihrer Wortart, sondern nach ihrer \textbf{Bedeutung} und \textbf{Funktion} ordnet, dann gehören sie zum Bereich der \textbf{kausalen Ausdrücke}.
    
    Man kann dabei zwei Hauptgruppen unterscheiden.
    
    \subsection{Ausdrücke des Grundes oder der Ursache}
        Zu dieser Gruppe gehören Mittel, die den Anlass oder die Ursache nennen.
        
        \paragraph{Typische Beispiele}
            \textbf{wegen}, \textbf{aufgrund}, \textbf{weil}, \textbf{da}, \textbf{aus}, \textbf{vor}, \textbf{durch}
    
    \subsection{Ausdrücke der Folge oder Konsequenz}
    Zu dieser Gruppe gehören Mittel, die eine Konsequenz aus einem vorher genannten Grund ausdrücken.
    
    \paragraph{Typische Beispiele}
        \textbf{deshalb}, \textbf{deswegen}, \textbf{darum}, \textbf{daher}, \textbf{aus diesem Grund}
    
    
    \section{Zusammenfassung}
    Kausale Mittel sind sprachliche Werkzeuge, mit denen man im Deutschen Ursache-Folge-Beziehungen ausdrückt. Manche nennen direkt den \textbf{Grund}, andere die \textbf{Folge}. Besonders wichtig ist die Unterscheidung zwischen Formen, die einen Grund einführen, und Formen, die eine Konsequenz formulieren.
    
    \textbf{wegen} ist ein Ausdruck des \textbf{Grundes}.\\
    \textbf{weil} und \textbf{da} sind ebenfalls typische Mittel für den \textbf{Grund}.\\
    \textbf{deshalb}, \textbf{deswegen}, \textbf{darum} und \textbf{daher} sind typische Ausdrücke der \textbf{Folge}.\\
    Zusammen gehören sie semantisch zum Bereich der \textbf{kausalen Ausdrücke} oder des \textbf{kausalen Zusammenhangs}.

%\bibliographystyle{plainnat}
%\bibliography{/home/donkarlo/Dropbox/repo/nd_shared_research/bibliography/bibliography.bib}

\end{document}